\section{Introduction}
\label{sec:intro}
A multicriteria analysis often produces several efficient alternatives while a decision maker ultimately needs a recommendation. Selecting one alternative requires assumptions about acceptable trade-offs, even when those assumptions are expressed without an explicit criterion-weight vector. Weighted sums, reference-point methods and distance-based methods provide different ways of making these assumptions operational \citep{hwang1981,wierzbicki1980,miettinen1999}. Their recommendations also depend on the numerical scales on which criteria are compared. When normalisation bounds are estimated, robustness to those bounds is a distinct concern from uncertainty about preferences.

This paper considers a declared family of monotone questions about each alternative, called probes. A probe can assess one criterion, the mean across criteria, or the worst criterion. Its value is transformed into an active-range disappointment: zero denotes the best probe value available in the fixed candidate set and one the worst. A decision maker can inspect the named concerns before accepting the recommendation. The probe family, its multiplicities, the candidate set and the normalisation convention remain substantive modelling choices.

Exact lexicographic minimax orders the sorted disappointments from worst to best. It gives any improvement in the largest disappointment priority over all improvements below it, regardless of the size of that leading difference. This is appropriate when priorities are strict. It can be difficult to justify when the decision model has a finite reporting resolution. Pairwise difference tolerances are an imperfect remedy because approximate equality need not be transitive. We instead use fixed resolution categories, followed by exact refinements, so that the comparison is a lexicographic order on a single well-defined vector.

The resulting rule, ReLexTail, first compares categories of the maximum and of upper-tail mean disappointments. Only when all category coordinates agree does it compare exact tail means and then the full sorted profile. Its contribution is this explicit hierarchy together with its decision-aiding interpretation, finite-set properties and treatment of uncertain normalisation bounds. The building blocks, including ordered aggregation, empirical conditional value-at-risk (CVaR) and lexicographic optimisation, are established tools. We do not claim their invention or a general superiority over other MCDA approaches.

\begin{figure}[tbp]\centering
\includegraphics[width=\linewidth]{pipeline.pdf}
\caption{Overview of the decision procedure and graphical abstract. The declared model determines active-range disappointments. Category comparisons precede exact refinement. Under uncertain bounds, sampling and interval certification answer different questions and require different reports.}
\label{fig:pipeline}\end{figure}

Ordered weighted averaging assigns weights to ordered outcomes rather than to named criteria \citep{yager1988}. Linear formulations of ordered aggregation and sums of the largest functions provide useful computational machinery \citep{ogryczak2003owa,ogryczak2003ksum}. Equitable optimisation studies related order-based preferences \citep{kostreva1999}. ReLexTail uses ordered quantities to evaluate a declared multiset of concerns. Its categories postpone sub-resolution differences, but its exact refinements still distinguish candidates within a category-optimal set. Accordingly, the categories do not create permanent indifference in the final preorder.

CVaR provides a continuous summary of an upper tail, including fractional empirical atoms \citep{rockafellar2000}. Here its atoms are probe disappointments, not uncertain future outcomes. A separate CVaR is used only in experimental evaluation, across held-out preference vectors. Keeping these two distributions distinct is essential: a small tail of declared concerns need not imply a small tail of regret for an external decision maker.

Stochastic multicriteria acceptability analysis explores the support for alternatives under uncertain preference parameters and data \citep{lahdelma1998,tervonen2007}. Robust ordinal regression considers compatible value functions and distinguishes necessary from possible preferences \citep{greco2008}. Our possible and necessary winner sets use similar quantifiers but vary normalisation bounds while holding the declared preference rule fixed. They therefore complement preference-uncertainty analysis rather than replace it. Likewise, exact LexPR in this paper denotes a reference order on the same disappointment representation. It is a baseline definition, not a claim that lexicographic minimax itself is new.

Recent work reinforces four adjacent needs without making them interchangeable. Robustness coefficients diagnose ranking sensitivity \citep{paradowski2025coefficients}; unknown-weight procedures explore preference uncertainty \citep{wieckowski2025unknown}; Bayesian pairwise queries support interactive elicitation \citep{huber2025bayesian}; and contrastive records can explain why one alternative is preferred to another \citep{erwig2025explaining}. Distributionally robust optimisation instead protects against ambiguity in a probability law \citep{kuhn2025dro}. ReLexTail addresses none of these problems in full generality: its primary uncertainty is a bounded perturbation of normalisation parameters, and its explanation is a deterministic replay of the declared profile.

\begin{table}[tbp]\centering\small
\caption{Structural positioning of neighbouring approaches. Entries describe the primary role of each family, not a quality ranking; implementations within a family can add further capabilities.}
\label{tab:positioning}
\begin{tabular}{@{}>{\raggedright\arraybackslash}p{.17\linewidth}>{\raggedright\arraybackslash}p{.29\linewidth}>{\raggedright\arraybackslash}p{.23\linewidth}>{\raggedright\arraybackslash}p{.20\linewidth}@{}}\toprule
Approach & Primary modelling object & Primary uncertainty & Typical output\\\midrule
Weighted score/distance & Compensatory score or distance & Weights/data if varied & Score or ranking\\
Leximax / LexPR & Ordered worst-to-best profile & None in the order & Exact optimal class\\
SMAA & Distribution over preferences/data & Preference/data sampling & Acceptability indices\\
Robust ordinal regression & Compatible preference models & Preference imprecision & Necessary/possible relations\\
DRO & Worst distribution in an ambiguity set & Probability-law ambiguity & Robust decision\\
ReLexTail & Fixed cells, tail summaries, exact profile & Bounded normalisation parameters & Category/exact winner sets\\\bottomrule
\end{tabular}
\end{table}

Certifying a recommendation under uncertain bounds raises a question that is easy to answer too conservatively. The natural sufficient condition asks every candidate to keep every risk coordinate inside its resolution cell, and it fails whenever any coordinate touches a cell boundary -- including a coordinate belonging to an alternative that is not in contention, a coordinate that is structurally invariant, and a coordinate occurring after the comparison has already been decided. None of those motions can change the recommendation. A second, independent source of conservatism sits in the arithmetic: the two normalisation anchors of a probe are shared by all candidates at a given bound vector, and widening the numerator and denominator of the disappointment ratio into separate intervals lets each take an extreme value the shared anchor could not have produced. Section~\ref{sec:certrelex} replaces both, proves that neither replacement weakens soundness, and measures what each is worth on a prespecified benchmark with locked splits.

The contribution is eightfold. First, we specify a fixed-cell tail-regret preorder with exact refinement, including boundary and tie conventions. Second, we establish finite-set existence, transitivity, conditional Pareto compatibility, replication invariance and a category-stability condition. Third, we distinguish sampled sensitivity from sound interval possible-winner elimination. Fourth, we give a sequential polyhedral formulation and its stage count. Fifth, we prove a decision-focused certification condition that requires no category to be retained, and exhibit a rational instance on which the all-candidate condition certifies nothing while the decision-focused condition certifies the entire uncertainty box without subdivision. Sixth, we prove that keeping the shared probe anchors inside a single quotient yields an enclosure contained in -- and generically strictly inside -- the independent one, and exhibit a rational instance where that difference decides whether the box is certified at all. Seventh, we test the quality--stability trade-off with matched perturbations, held-out preferences, paired uncertainty summaries, grid-origin sensitivity and ablations on 90 synthetic sets, plus a public building-simulation demonstration, and we test the certification claims on a separate locked benchmark with an exact rational oracle and an independent replay verifier. Eighth, we produce contrastive decision records and a complete replay package. These contributions combine established ordered-risk tools around finite decision resolution; they do not imply universal superiority or operational validation. The certification results say nothing about decision quality, external regret or calibration, which are not tested here. A ten-supplier illustration supplies an interpretable worked decision.
